\section{PUBLICAÇÕES}

Parte dos resultados desenvolvidos nesta dissertação resultou em produção bibliográfica diretamente relacionada às contribuições metodológicas apresentadas. O método \gls{CLAIRE} \cite{ferreira2025claire} foi publicado em periódico científico \textit{\textbf{Machine Learning}} da \textit{\textbf{Springer Nature}}, no qual é proposta uma abordagem probabilística para avaliação de métodos de clustering baseada na \gls{TRI}, modelando a concordância entre modelos como matriz de respostas e estimando habilidades latentes sem acesso a rótulos verdadeiros.

Além disso, o modelo $\beta^4$-\gls{IRT} \cite{ferreirajunior2023beta4irtnewbeta3irtenhanced}, que introduz uma parametrização para separar sinal e magnitude do parâmetro de discriminação, mitigando problemas de não-identificabilidade observados no $\beta^3$-\gls{IRT}, foi disponibilizado como preprint na plataforma \textbf{arXiv}. Nesse trabalho, são investigadas as propriedades teóricas do modelo e sua capacidade de recuperação de parâmetros em cenários com instâncias ruidosas ou ambíguas no contexto de Aprendizado de Máquina.
